\begin{figure}[tbp]
\centering
\begin{figurealt}{GLTF or GLB enters vendored cgltf for parsing and buffers. CNA's shared import core validates and extracts scene nodes, mesh groups, materials, skins, animation records, cameras, and diagnostics. The runtime front end converts all represented groups into live GraphicsDevice resources in one Model. The offline gltf to cnj front end writes one CNJ descriptor and packed sidecar set per group. Runtime and offline share parsing and extraction but not final object or file construction.}
\resizebox{\textwidth}{!}{%
\begin{tikzpicture}[node distance=9mm and 12mm]
  \node[cna public] (input) {\texttt{.gltf} / \texttt{.glb}};
  \node[cna internal, right=of input] (cgltf) {cgltf parse +\\buffer/image access};
  \node[cna internal, right=of cgltf, minimum width=36mm] (core) {shared import core\\scene, groups, materials,\\skins, animation, report};
  \node[cna internal, above right=of core] (runtime) {runtime reader\\all represented groups};
  \node[cna public, right=of runtime] (model) {live buffers, effects,\\textures, \texttt{Model}};
  \node[cna internal, below right=of core] (offline) {\texttt{gltf\_to\_cnj}\\offline front end};
  \node[cna public, right=of offline] (sidecars) {per-group CNJ +\\packed sidecars};
  \node[cna warning, below=of core] (boundary) {parse $\ne$ runtime use $\ne$\\renderer draw $\ne$ pixel proof};
  \draw[cna flow] (input) -- (cgltf);
  \draw[cna flow] (cgltf) -- (core);
  \draw[cna flow] (core) -- (runtime);
  \draw[cna flow] (runtime) -- (model);
  \draw[cna flow] (core) -- (offline);
  \draw[cna flow] (offline) -- (sidecars);
  \draw[cna optional] (boundary) -- (core);
\end{tikzpicture}}
\end{figurealt}
\caption{The shared glTF import core and its two front ends at \CnaRevisionShort. Shared parsing
does not make the runtime and offline output routes equivalent.}
\label{fig:gltf-import-routes}
\end{figure}
